\smallskip
\noindent\textbf{Theoretical analysis of the bidding schemes.}
We now show that both \poorman and \allpay implement the reward-aware urgency ordering whenever one goal has a sufficiently large urgency advantage over the others.
To capture stochastic motion of goals, we formulate the urgency scores using conditional expectations under the transition kernel of the \momdp.
The deterministic stationary-goal case discussed above is then recovered as a special case in which these conditional expectations collapse to point values.
\KM{No need to talk about the deterministic case here.}
\KM{Update the text}

Suppose $\pi_i$ is the local policy for the $i$-th goal, mapping every state $s\in S$ and time $t$ to a probability distribution over the bid $b_i$ and a deterministic action $a_i$.
For convenience, we will write $\pi_i^b\colon s\mapsto \delta_i([0,\beta])$ and $\pi_i^a\colon s\mapsto a_i$ as two separate components of the policy, where $\delta_i([0,\beta])$ represents a probability distribution over $[0,\beta]$, and $b_i\sim \delta_i([0,\beta])$ will represent a concrete bid value sampled from this distribution.
Given the state $s$, and given fixed $\pi_i^a$ and $\pi_i^b$ for all $i\in [1;m]$, we obtain a Markov chain starting at $s$, which is an \momdp with a singleton action set, with the deterministic initial state $s$, i.e., $\mu(s)=1$.
This Markov chain induces a unique probability measure on the sample space of paths and rewards, and let $V_{\pi_i,\pi_{-i}}(s)$ denote the \textit{expected} discounted sum of reward obtained by the goal $i$ in this Markov chain.
(The notation $\pi_{-i}$ indicates the collection of every $j$-th policy for $j\neq i$.)
A policy profile $\{\pi^*_i\}_{i\in [1;m]}$ is called a \textit{Nash equilibrium in the subgame starting at} $s$ if for every $i$:
$	
\forall \pi_i\neq \pi^*_i\;.\;V_{\pi^*_i,\pi^*_{-i}}(s) \geq V_{\pi_i,\pi^*_{-i}}(s).
$
To simplify the notation, we will write $V_i^*(s)$ to denote the value $V_{\pi^*_i,\pi^*_{-i}}(s)$.

We will now perform a backward induction over time to obtain the subgame perfect equilibrium, i.e., the equilibrium strategy profile that is a Nash equilibrium in every subgame.
Suppose we are currently at the state $s$, and let $i$ be a given policy index. 
We consider two cases:
(a) $i$ wins the bidding, in which case $i$ picks $\pi_i^a$ to maximize its own reward, and the optimal reward for $i$ is: $\widetilde V_i^{\win}(s) \coloneqq \max_{\pi_i^a} \mathbb{E}[\gamma^\tau V_i^*(s^\tau)+\sum_{t=0}^{\tau-1} \gamma^t r_i(s^t,a^t)]$, where $s^0=s$ and $s^0a^0\ldots s^\tau a^\tau$ is the random trace obtained by executing $\pi_i^a$ successively for $\tau$ steps (no bidding takes place during this time);
(b) $i$ loses the bidding, in which case another policy $j$ picks $\pi_j^a$ to maximize its own reward, and by taking the minimum over all $j$, we obtain the worst-case resulting reward for $i$:
$\widetilde V_i^{\lose}(s) \coloneqq \min_{j\neq i,\pi_j^{a*}} \mathbb{E}[\gamma^\tau V_i^*(s^\tau) + \sum_{t=0}^{\tau-1} \gamma^t r_i(s^t,a^t)]$, where $\pi_j^{a*}$ is the policy that achieves $\widetilde V_j^{\win}(s)$, $s^0=s$, and $s^0a^0\ldots s^\tau a^\tau$ is the random trace obtained by executing $\pi_i^a$ successively for $\tau$ steps.
Intuitively, $\widetilde V_i^{\win}(s)$ and $\widetilde V_i^{\lose}(s)$ are the expected rewards before the payment of the bid value, and if the concrete bid value is $b_i$, then after adjusting the bid payment, the net payoffs become:
\begin{align*}
	V_i^{\win}(s,b_i) &\coloneqq \widetilde V_i^{\win}(s) - \rho b_i,\\
	V_i^{\lose}(s,b_i) &\coloneqq \widetilde V_i^{\lose}(s) - \rho b_i\cdot \mathbf{1}[\text{\allpay}].
\end{align*}
Given a set of fixed bidding policies $\{\pi_i^b\}_{i\in [1;m]}$, and for the set of sampled bids $\{b_i\sim\pi_i^b(s)\}_{i\in [1;m]}$, define:
\begin{align*}
	U_i(s,b_i,b_{-i}) \coloneqq
		\begin{cases}
			V_i^{\win}(s,b_i) & \text{if } b_i > \max_{j\neq i} b_j,\\
			V_i^{\lose}(s,b_i) & \text{if } b_i < \max_{j\neq i} b_j,\\
			\frac{1}{|W|}V_i^{\win}(s,b_i) + \left(1-\frac{1}{|W|}\right)V_i^{\lose}(s,b_i) & \text{if } i\in W,
		\end{cases}
\end{align*}
with $W\coloneqq \{k\in [1;m]\mid b_k=\max_{\ell\in [1;m]} b_\ell\}$ being the set of highest bidders.
Suppose $\{\pi_i^{b*}\}_{i\in [1;m]}$ is a Nash equilibrium for the set of utilities $\{U_i(s,\cdot,\cdot)\}_{i\in [1;m]}$, i.e., for every $i$:
\begin{align*}
	\forall \pi_i^b\neq \pi_i^{b*}\;.\;\mathbb{E}_{b_i\sim \pi_i^{b*},b_{-i}\sim\pi_{-i}^{b*}}[U_i(s,b_i,b_{-i})] \geq \mathbb{E}_{b_i\sim \pi_i^{b},b_{-i}\sim\pi_{-i}^{b*}}[U_i(s,b_i,b_{-i})].
\end{align*}
Then it follows from the backward induction that:
\begin{align*}
	V_i^*(s) = \mathbb{E}_{b_i\sim \pi_i^{b*},b_{-i}\sim\pi_{-i}^{b*}}[U_i(s,b_i,b_{-i})].
\end{align*}
The \emph{net gain} for bidding $b_i$ at $(s,t)$ is defined as:\KM{$\Delta$ is used before for probability distributions.}
\begin{align*}
	\Delta_i^{\mathrm{net}}(s,b_i)&\coloneqq V^{\win}_i(s,b_i) - V^{\lose}_i(s,b_i)\\
	&=  \Delta^{\mathrm{gross}}_i(s) - \rho b_i\cdot \mathbf{1}[\text{\poorman}],
\end{align*}
where the \emph{gross gain} is
\begin{align*}
	\Delta^{\mathrm{gross}}_i(s)\coloneqq \widetilde V_i^{\win}(s)-\widetilde V_i^{\lose}(s).
\end{align*}
The quantity $\Delta^{\mathrm{gross}}_i(s,t)$ is exactly the reward-aware urgency of policy $i$: it measures the continuation-value loss incurred by not allocating the next bidding interval to $i$.
\GS{Talk about how this urgency measure, in addition to the deadline, also encodes cooperation between policies, if there is any, because if \(V_i^{\lose}(s,t)\) is high, i.e., then the urgency measure is low.}
Thus RAUF selects a policy maximizing $\Delta^{\mathrm{gross}}_i(s,t)$.

The following theorem establishes that the \poorman setting implements RAUF whenever the leading urgency exceeds the runner-up by more than one bid increment.

\begin{theorem}
	In the \poorman setting, consider a single bidding round at the agent state $s\in S$.
 	Suppose the parameters $\rho$ and $\beta$ are so chosen that $\rho\beta > \max_{j\in [1;m]}\Delta_j^{\mathrm{gross}}(s)$.
	Let $i^*$ be the goal index with the largest gross gain, i.e.,
	\begin{align}\label{eq:all pay:gross gain gap}
		\Delta_{i^*}^{\mathrm{gross}}(s) > \max_{j\neq i^*}\Delta_j^{\mathrm{gross}}(s).
	\end{align}
	Then in any pure Nash equilibrium, policy $i^*$ wins control.
\end{theorem}

\begin{proof}
	Suppose for contradiction, some other policy $j\neq i^*$ wins control by bidding $b_j$, while policy $i^*$ bids $b_{i^*}<b_j$. 
	As this is a Nash equilibrium, it must hold that 
	$\Delta_j^{\mathrm{net}}(s,b_j) > 0$, i.e., $ \Delta_j^{\mathrm{gross}}(s) - \rho b_j > 0$,
	because otherwise, policy $j$ would be better off by losing the bidding instead.
	This gives us the upper bound $\rho b_j < \Delta_j^{\mathrm{gross}}(s)$.
	It follows from $\rho\beta > \max_{k\in [1;m]}\Delta_k^{\mathrm{gross}}(s)$ that $\rho b_j < \rho\beta$, and hence $b_j<\beta$, i.e., the policy $j$ does not make the highest possible bid.
	
	Now we assumed that policy $i^*$ had bid $b_{i^*}<b_j$, but we show that policy $i^*$ would achieve a higher net gain by bidding $b_j+\varepsilon$ for any $0<\varepsilon< \min\{ (\Delta^{\mathrm{gross}}_{i^*}(s) - \Delta_j^{\mathrm{gross}}(s))/\rho, \beta-b_j\}$, contradicting the Nash equilibrium:	 
$	\Delta_{i^*}^{\mathrm{net}}(s,b_j+\varepsilon) 
	= \Delta^{\mathrm{gross}}_{i^*}(s) - \rho(b_j+\varepsilon) 
	= \Delta^{\mathrm{gross}}_{i^*}(s) - \rho b_j -\rho\varepsilon  
	> \Delta^{\mathrm{gross}}_{i^*}(s) - \Delta_j^{\mathrm{gross}}(s) -\rho\varepsilon
	> 0.
$
	Here, notice that $b_j+\varepsilon < \beta$, which is a valid bid.
\end{proof}


The following theorem establishes that the \allpay setting concentrates probability mass on the RAUF choice whenever the top urgency is strictly ordered above the remaining urgencies.

\begin{theorem}
	In the \allpay setting with $m\geq 3$, consider a single bidding round at the agent state $s\in S$ and time $t$.
	Let $i^*$ and $j^*$ be policies for which the following gross gain order holds:
	\begin{align}\label{eqn:allpay urgency gap}
		\Delta_{i^*}^{\mathrm{gross}} > \Delta_{j^*}^{\mathrm{gross}} > \max_{k\notin \{i^*,j^*\}}\Delta_k^{\mathrm{gross}}.
	\end{align}
	Then there is no Nash equilibrium with pure strategies, although a unique Nash equilibrium exists with mixed strategies (where at each $(s,t)$, each policy uses a probability distribution to select the bid).
	In this Nash equilibrium, policy $i^*$ wins control with probability $1-\sfrac{\Delta_{j^*}^{\mathrm{gross}}}{2\Delta_{i^*}^{\mathrm{gross}}}$.	
\end{theorem}

\begin{proof}
	It follows from the work of Hillman and Riley~\cite{hillman1989politically} that in the Nash equilibrium, the optimal bid of the policies are as follows.
	The policy $i^*$ picks a uniformly random bid $b_i^*$ from $[0,\Delta_{j^*}^{\mathrm{gross}}/\rho]$, the policy $j^*$ picks a uniformly random bid $b_j^*$ from $(0,\Delta_{j^*}^{\mathrm{gross}}/\rho]$ while picking the bid $0$ with probability $(\Delta_{i^*}^{\mathrm{gross}} - \Delta_{j^*}^{\mathrm{gross}})/\Delta_{i^*}^{\mathrm{gross}}$, and every other policy deterministically picks the bid $0$.
	Clearly, if $b_i^*>b_j^*$, then $i^*$ wins control, and we need to show that $P[b_i^* > b_j^*]  = 1- \sfrac{\Delta_{j^*}^{\mathrm{gross}}}{2\Delta_{i^*}^{\mathrm{gross}}}$.
	To show this, consider the probability density function $f_{i^*}(b) = \sfrac{\rho }{\Delta_{j^*}^{\mathrm{gross}}}$ of $b_{i^*}$, and the cumulative distribution function $F_{j^*}(b) = (\Delta_{i^*}^{\mathrm{gross}} - \Delta_{j^*}^{\mathrm{gross}})/\Delta_{i^*}^{\mathrm{gross}} + \sfrac{\rho b}{\Delta_{i^*}^{\mathrm{gross}}}$ of $b_{j^*}$.
	Then,
	\begin{align*}
		P[b_i^* > b_j^*] 
		&= \int_{-\infty}^\infty F_{j^*}(b)f_{i^*}(b)db\\
		&= \int_{0}^{\Delta_{j^*}^{\mathrm{gross}}/\rho} \left(\frac{\Delta_{i^*}^{\mathrm{gross}} - \Delta_{j^*}^{\mathrm{gross}}}{\Delta_{i^*}^{\mathrm{gross}}} + \frac{\rho b}{\Delta_{i^*}^{\mathrm{gross}}}\right)\cdot \frac{\rho }{\Delta_{j^*}^{\mathrm{gross}}} \cdot db \\
		&= \int_{0}^{\Delta_{j^*}^{\mathrm{gross}}/\rho} \frac{\Delta_{i^*}^{\mathrm{gross}} - \Delta_{j^*}^{\mathrm{gross}}}{\Delta_{i^*}^{\mathrm{gross}}} \cdot \frac{\rho }{\Delta_{j^*}^{\mathrm{gross}}} \cdot db
		+ \int_{0}^{\Delta_{j^*}^{\mathrm{gross}}/\rho} \frac{\rho b}{\Delta_{i^*}^{\mathrm{gross}}}\frac{\rho }{\Delta_{j^*}^{\mathrm{gross}}} \cdot db\\
		&=  \frac{\Delta_{i^*}^{\mathrm{gross}} - \Delta_{j^*}^{\mathrm{gross}}}{\Delta_{i^*}^{\mathrm{gross}}} \cdot \frac{\rho }{\Delta_{j^*}^{\mathrm{gross}}} \cdot \frac{\Delta_{j^*}^{\mathrm{gross}}}{\rho}
		+ \frac{\rho^2}{\Delta_{i^*}^{\mathrm{gross}}\Delta_{j^*}^{\mathrm{gross}}}\cdot \frac{\left(\Delta_{j^*}^{\mathrm{gross}}\right)^2}{2\rho^2}\\
		&= \frac{\Delta_{i^*}^{\mathrm{gross}} - \Delta_{j^*}^{\mathrm{gross}}}{\Delta_{i^*}^{\mathrm{gross}}} + \frac{\Delta_{j^*}^{\mathrm{gross}}}{2\Delta_{i^*}^{\mathrm{gross}}}\\
		&= 1-\frac{\Delta_{j^*}^{\mathrm{gross}}}{\Delta_{i^*}^{\mathrm{gross}}} + \frac{\Delta_{j^*}^{\mathrm{gross}}}{2\Delta_{i^*}^{\mathrm{gross}}}\\
		&= 1-\frac{\Delta_{j^*}^{\mathrm{gross}}}{2\Delta_{i^*}^{\mathrm{gross}}}
	\end{align*}
\end{proof}

If the second inequality in \eqref{eqn:allpay urgency gap} becomes an equality, then it is known that the Nash equilibrium is no longer unique, and in fact, there will be a continuoum of Nash equilibria~\cite{baye1996all}, and reaching the desirable equilibrium will require separate mechanisms.
Instead, we argue that the event of any of the inequalities in \eqref{eqn:allpay urgency gap} becoming an equality is a measure zero event, and therefore it is vanishingly unlikely in practice.
Moreover, when $\Delta_{i^*}^{\mathrm{gross}}$ is substantially larger than $\Delta_{j^*}^{\mathrm{gross}}$, the winning probability
\[
1-\frac{\Delta_{j^*}^{\mathrm{gross}}}{2\Delta_{i^*}^{\mathrm{gross}}}
\]
approaches $1$.
Hence, although \allpay need not deterministically implement RAUF, it still concentrates control on the most reward-aware urgent policy whenever the top-two urgency gap is large.

